% =====================================================================================
%  Chapter 5 — Experimental Setup and Protocol
%  Thesis: Comparison of Control Strategies for Interaction-Force Aware Multirotor Teams
%
%  DRAFT, 23 August 2026 (writing-track item W.7).
%  REVISED 2026-10-14 (thesis-revision run, Phase 4). Markers: % NEW, % CHANGED, % DECISION,
%  % VERIFY, % REMOVED (H1: original text kept as a comment, not printed). Plan rows P5-xx refer to
%  docs/thesis_revision/stage3_revision_plan.md.
%  Sources: docs/05_Experimental_Protocol_2Robot.md, docs/10_Formation_Library.md,
%           docs/11_Hardware_Readiness_Checklist.md, docs/09_Simulation.md.
%           Design facts re-checked against firmware_app/src/lib.rs, firmware_app/traj_iface.c,
%           tools/usd_study_config.txt and crazyswarm2/.../formations/scenarios.py.
%
%  HONESTY NOTE (updated 21 September 2026). Hardware flights began September 2026 (C.0/C.4
%  indicative runs, uSD pipeline). C.1 collection and full C.4 repeats were still in progress
%  when this draft was last edited. Revise tense section-by-section as the campaign completes;
%  record protocol deviations here, not only in the results chapters.
%
%  Citation baseline: 2026-08-27 (see docs/thesis/CITATION_BASELINE.md).
%  After that date: only NEW cites or EDITED prior-work sentences need re-checking.
% =====================================================================================

\chapter{Experimental Setup and Protocol}
\label{ch:setup}

Chapter~\ref{ch:control} described Strategies~0--3, each as its own controller. This chapter
describes the platform, scenarios, measurement chain and protocol. The design removes alternative
explanations for an observed difference.

Section~\ref{sec:setup-validity} states the threats this design is built against, and which of
them it does not eliminate.

% DECISION D-05 (H9, lab status): the following paragraph reports campaign progress, vehicle
%   nicknames, dates, a manifest file name and an internal document number. Options: (i) delete
%   (recommended); (ii) replace by a future-tense design statement ("The campaign will ..."); (iii)
%   neutral implementation-status wording without flight claims. Text left unchanged.
\paragraph{Campaign status (September 2026).}
Two-robot flights with OptiTrack and onboard logging at \SI{500}{\hertz} have been executed on
the brushless pair \textbf{cf5} (study vehicle) and \textbf{cf\_second} (top). Micro-SD merges
with per-vehicle role alignment replaced radio-only windows for thesis metrics. On
\textbf{2026-09-21}, eight verified geometric merges were archived (five C.1 training-eligible
A1/A3 cells plus three A8 QA shakedowns; manifest
\texttt{c1\_2026\_09\_21\_merged/manifest\_2026\_09\_21\_c1.json}), with dashboards and phase
metrics produced by the standard analysis suite. Indicative Strategy~0 vs~1 comparisons exist from
earlier sessions; the systematic comparison (C.4) still requires repeated flights per controller
and scenario. Residual-model training data collection (C.1) \emph{has begun} under
\textbf{geometric control on cf5} so measured $a_{\mathrm{res}}$ is not cancelled by full INDI,
but the planned dz/speed grid is only partly filled (see \texttt{docs/25} coverage table).

% -------------------------------------------------------------------------------------
\section{Platform}
\label{sec:setup-platform}

% CHANGED (V4): the motor spacing from the firmware platform definition. crazyflie-firmware
%   src/platform/interface/platform_defaults_cf21bl.h:46 (ARM_LENGTH 0.050 m, centre to motor);
%   lib.rs:261 (ARM_M = sqrt(2)/2 * 0.050). Diagonally opposite motors are therefore 0.10 m apart.
%   Mass: 41 g is the configured mass of the study vehicles (crazyflies.yaml, mass 0.041); the
%   firmware default is 39.3 g for a different deck set (platform_defaults_cf21bl.h:47-52).
The vehicles are Crazyflie~2.1 brushless quadrotors, nominally \SI{41}{\gram}, with
\SI{10}{\centi\metre} between diagonally opposite motors. All vehicles in a team are the same platform: the gain block
is shared across robots, so a mixed team would fly at least one vehicle on the wrong mass, thrust
constant and inertia, which would appear in the logs as a residual and be indistinguishable from
an interaction force (Section~\ref{sec:model-measure-fail}).

Two vehicles are the minimum for the protocol. A third is required for the three-robot scenarios,
which cannot be run without it.

Each vehicle carries two additions that are not optional for this work:

% CHANGED (P5-20): the default rotor-speed source.
\paragraph{A rotor-speed source}, normally the rotor speed that the electronic speed controllers
report over the bidirectional DShot protocol, or alternatively an optical rotor-speed (RPM) deck;
a parameter selects which one the controller reads. Without it the applied
thrust cannot be reconstructed, so $a_{\mathrm{res}}$ is unavailable and \emph{reads identically
zero} rather than failing loudly. Strategy~1 additionally cannot run at all. This is the
single most consequential hardware prerequisite in the setup.
% CHANGED: code, traj_iface.c g_indi_rpm_source (default 1 = DShot, L508-520) and rpm_get_all
%   (L701-730).

\paragraph{A micro SD (uSD) deck}, the recording medium for the dataset, for the reasons in
Section~\ref{sec:setup-logging}.

% CHANGED (P5-01): the control loop runs at 1 kHz.
State estimation uses the on-board extended Kalman filter, fed with external pose from a
motion-capture (mocap) system at \SI{100}{\hertz}. The control loop runs onboard at
\SI{1000}{\hertz}; the dataset is recorded at \SI{500}{\hertz}, and the learned model of
Strategy~2 is evaluated at about \SI{100}{\hertz} (Chapter~\ref{ch:model}).
% CHANGED: code, crazyflie-firmware stabilizer.c (RATE_MAIN_LOOP = 1000 Hz, controller called every
%   step); lib.rs:controllerOutOfTree (dt from 1 ms ticks, L2350-2351); usd_study_config.txt (500 Hz);
%   traj_iface.c g_rnn_div = 10.

% -------------------------------------------------------------------------------------
\section{Flight volume}
\label{sec:setup-volume}

The tracked volume measures approximately \SI{2}{\metre} $\times$ \SI{4}{\metre} horizontally and
extends to \SI{1.7}{\metre} in height. Two distinct limits are enforced and should not be
conflated: the \emph{tracked volume}, outside which the motion-capture system cannot see the
vehicles, and an \emph{operational floor} of \SI{0.30}{\metre}, chosen so that ground effect does
not contaminate the measurement of inter-vehicle interaction.

The floor is lowered to \SI{0.10}{\metre} for one scenario only, the near-ground pass of
Section~\ref{sec:setup-scenarios}, whose purpose is to measure ground effect in isolation so that
it can be separated from downwash rather than confounded with it.

The volume is a binding constraint rather than a formality. Several scenarios sit within
approximately \SI{10}{\centi\metre} of a wall at their extremes, and the scenario definitions are
automatically centred and, where necessary, rotated to fit. That a configuration fits is checked
before flight, not discovered during it.

\todo{The volume figures are corroborated from two independent statements but have not been
tape-measured. Measure them and update this section.}

% -------------------------------------------------------------------------------------
\section{The formation library}
\label{sec:setup-scenarios}

% CHANGED (P5-13): the table is referenced, and parameters are given.
Sixteen scenarios, fixed before any data collection and not modified afterwards, are listed in
Table~\ref{tab:setup-scenarios} with their default parameters. Freezing them is
what allows every strategy to meet identical excitation; a library that evolved during the campaign
would make later strategies incomparable with earlier ones.

The scenarios are parameterised rather than hard-coded, separations, offsets and speeds are
arguments, so that a separation sweep is the same scenario at different parameters rather than
a different scenario. Each is compiled to a per-vehicle polynomial trajectory and uploaded through
the standard trajectory interface, so that the same command produces the same motion in simulation
and on hardware. The commanded relative geometry between any two vehicles is known in advance from
the scenario definition, which the analysis uses to define its windows
(Section~\ref{sec:setup-metrics}).

% CHANGED (P5-13): vehicle count and default parameters per scenario, from scenarios.py.
\begin{table}[t]
  \centering
  \caption{The formation library. $n$ is the default number of vehicles; separations are the
    defaults of the scenario definitions and can be varied as parameters. Scenarios marked $\dagger$
    carry elevated risk and are gated behind an explicit acknowledgement; those marked $\ast$ are
    controls with no vertical overlap, in which no downwash is expected (C2 also probes the wake
    behind a vehicle).}
  \label{tab:setup-scenarios}
  \small
  \begin{tabular}{@{}llcp{0.25\linewidth}p{0.36\linewidth}@{}}
    \toprule
    & Scenario & $n$ & Default parameters & Purpose \\
    \midrule
    \multicolumn{5}{@{}l}{\emph{Two vehicles}}\\
    A1 & Vertical stack, hover      & 2 & $\Delta z = \SI{0.50}{\metre}$, \SI{12}{\second} hold & Baseline. Nothing moves, so any relative displacement is interaction rather than tracking error \\
    A2 & Vertical stack, tracking   & 2 & $\Delta z = \SI{0.30}{\metre}$, circle, \SI{0.4}{\metre\per\second} & Constant separation while both vehicles move \\
    A3 & Static top                 & 2 & $\Delta z = \SI{0.30}{\metre}$, \SI{0.4}{\metre\per\second}, 4 passes & Fixed wash source, lower vehicle sweeps through: captures the transition, not only the level \\
    A4 & Offset stack               & 2 & $\Delta z = \SI{0.60}{\metre}$, offset \SI{0.10}{\metre}, lemniscate & Partial overlap; asymmetric and hardest to model \\
    A5 & Reverse-circle stack       & 2 & $\Delta z = \SI{0.50}{\metre}$, radius \SI{0.40}{\metre}, \SI{7.5}{\second} period & Lateral offset swept twice per lap at constant height and speed \\
    A6 & Extreme stack $\dagger$    & 2 & $\Delta z = \SI{0.10}{\metre}$, \SI{0.20}{\metre\per\second} & Strongest interaction, smallest error budget \\
    A7 & Dynamic merge $\dagger$    & 2 & $\Delta z$ \SI{1.10}{\metre} $\to$ \SI{0.10}{\metre}, \SI{0.30}{\metre\per\second} & One continuous sweep of the whole separation range: gives a model the gradient, not only levels \\
    A8 & Vertical swap              & 2 & $\Delta z = \SI{0.25}{\metre}$, span \SI{1.0}{\metre}, \SI{6}{\second} & Wash arrives as a transient with a known arrival time; hardest case for a feedforward model \\
    \midrule
    \multicolumn{5}{@{}l}{\emph{Three vehicles}}\\
    B1 & I-stack                    & 3 & $\Delta z$ \SI{0.20}{\metre} and \SI{0.30}{\metre}, circle & Lowest vehicle in the combined wash of two \\
    B2 & V-stack                    & 3 & $\Delta z$ \SI{0.20}{\metre} and \SI{0.30}{\metre}, offset \SI{0.10}{\metre} & Asymmetric aggregation, which B1 cannot produce \\
    B3 & Vertical swap              & 3 & $\Delta z = \SI{0.22}{\metre}$, span \SI{0.55}{\metre} & Three crossings at three height differences in one flight \\
    \midrule
    \multicolumn{5}{@{}l}{\emph{Controls and coverage}}\\
    C1 & Side-by-side coplanar $\ast$ & 3 & spacing \SI{0.30}{\metre}, hover & A residual here, and not in A1, is not downwash \\
    C2 & Leader--follower line $\ast$ & 3 & gap \SI{0.50}{\metre}, \SI{0.4}{\metre\per\second} & The wake behind a moving vehicle, not the downwash below it \\
    C3 & Equilateral triangle $\ast$  & 3 & side \SI{0.50}{\metre}, hover & Coplanar three-body control \\
    C4 & Docking approach $\dagger$   & 2 & lateral \SI{1.0}{\metre}, $\Delta z$ \SI{0.50}{\metre} $\to$ \SI{0.10}{\metre} & Simultaneous lateral and vertical closure \\
    C5 & Near-ground pass & 1 & height \SI{0.15}{\metre}, \SI{0.25}{\metre\per\second} & Ground effect alone, separable from downwash \\
    \bottomrule
  \end{tabular}
\end{table}
% CHANGED: code, crazyswarm2/crazyflie_examples/crazyflie_examples/formations/scenarios.py
%   (builders A1-C5, L134-465; registry L466-470; gating: A6 extreme=True, A7/C4 extreme when the
%   final separation is below 0.15 m).

\subsection{The null scenarios}
\label{sec:setup-null}

C1--C3, the coplanar controls in which the vehicles fly side by side at equal height, are the most
important scenarios in the library and the least interesting to watch. They
place the vehicles in configurations where no significant inter-vehicle interaction should exist,
and they exist because of the property established in Section~\ref{sec:model-dynamics}: the
residual is a catch-all, non-zero even in isolated flight, so a residual measured in formation
cannot be attributed to the interaction without a reference.
% NEW: the physical basis, with its caveat, from a held source (R03 Sec. 5).
The physical basis for this choice is that a multirotor's downwash remains concentrated beneath
it; on their load stand, Gielis et al.\ found the aerodynamic influence of a vehicle negligible on
neighbours outside its downwash, even at small separations, while noting small, chaotic lateral
forces at larger distances~\cite{gielis2023aggregate}.

They serve two purposes. They bound the measurement floor, so that a reported improvement can be
distinguished from noise. And they detect model error: a residual that appears when no neighbour is
in the wash is a mis-identified mass or thrust constant, not an interaction, and must be corrected
before the campaign proceeds rather than explained afterwards.

\subsection{Regimes, and why the library spans them}

% CHANGED (P5-14): C2 moves; the others are static.
The scenarios are not a list of poses but a deliberate spread across the interaction regimes that
Research Question~2 distinguishes. A1, the two-robot vertical hover stack, and the coplanar
hover controls C1 and C3 are static, and the leader--follower control C2 translates at constant
speed; A2 and A3 are quasi-static; A4 and A5
sweep lateral offset; A7 sweeps separation continuously; A8 and B3 present the disturbance as a
transient with a known arrival time. A model fitted on one regime and evaluated on another is what
makes ``outside the training distribution'' a measurable condition rather than a caveat.

\begin{quote}
\textbf{Requirement for data collection: the set must excite lateral relative motion.}
The six network inputs are the relative position and the relative velocity of a neighbour, three
components each. A purely vertical scenario holds the four lateral components (relative $x$ and $y$,
position and velocity) nearly constant, so a model trained only on such data has seen almost no
variation in four of its six inputs and is extrapolating the first time a formation moves sideways. \textbf{The collection set must therefore include A4, A5
and A7 --- not only the vertical stacks A1, A2 and A3.} This is not a preference about coverage; a
set without them produces a model that cannot be evaluated on the scenarios it will meet.
\end{quote}

% NEW (P5-15): what the input gate implies for which scenarios exercise the network.
The input gate of the network (Section~\ref{sec:model-guards}) sharpens this requirement. A
neighbour more than \SI{0.2}{\metre} away along either horizontal axis contributes nothing to the
prediction, so the coplanar controls at their default spacings (\SI{0.30}{\metre} and wider) exercise
only the ground branch, A5 exercises the neighbour branch only during the part of each lap in which
the two vehicles pass within that band, and A4 at its default offset of \SI{0.10}{\metre} stays inside
it. Lateral excitation matters within the band, not across the whole room.

% REMOVED (P5-05, H1): "This was established in simulation before any flight: a dry run collecting
%   only a vertical stack produced training data in which two of the six network inputs never
%   varied, which the training pipeline reports as a degenerate normalisation statistic and which
%   no amount of training can repair." (a project simulation finding)
% NEW (P5-05): the requirement stated as a check that needs no flight.
The requirement follows from the inputs alone, and it can be checked before any flight from the
commanded trajectories of the planned collection set: an input whose commanded range is near zero
across the set will not vary in the data either.

% -------------------------------------------------------------------------------------
\section{Measurement and logging}
\label{sec:setup-logging}

\subsection{Why the dataset is recorded onboard}

% CHANGED (P5-07, H1-suspect): reasoning kept, project-measured numbers removed.
Radio telemetry shares a fixed bandwidth per radio dongle among all vehicles it serves, so the rate
at which each vehicle can stream its log falls as vehicles are added. Recording the dataset over
radio would therefore couple the sampling rate to the team size, which is unacceptable for a study
whose subject is what happens as vehicles are added.
% REMOVED (P5-07): "Radio telemetry saturates at approximately two vehicles: each drone transmits on
%   the order of 800 packets per second at full rate against a per-dongle ceiling near 1000."
% VERIFY: if the numbers are to be kept, cite the vendor's radio specification for the per-dongle
%   packet rate rather than the project's own observation.

% CHANGED (P5-02): 96 variables.
The micro SD deck records \num{96} variables at \SI{500}{\hertz} on each vehicle independently of
radio and of team size. Radio logging is retained for live monitoring and for single-vehicle
tuning, where the bandwidth exists.
% CHANGED: code, tools/usd_study_config.txt (header: on:fixedFrequency, 500; 96 variable lines).

\subsection{What is recorded}
\label{sec:setup-recorded}

% CHANGED (P5-03): the list completed from the logging configuration.
Position and velocity estimates and the attitude, as Euler angles and as a quaternion; gyro and
accelerometer; the measured residual $a_{\mathrm{res}}$; the attitude error $e_R$ and its norm; the
INDI torque and filtered angular acceleration; the network prediction $\hat{a}_{\mathrm{res}}$ and
its saturation flag; the commanded position, velocity and acceleration; the four motor commands;
the rotor speeds as reported by the motor controllers, as measured by the optical deck where
fitted, and as actually used by the controller after filtering; the battery voltage; the external
pose received from motion capture; and controller-internal diagnostic variables.
% CHANGED: code, tools/usd_study_config.txt (stateEstimate.*, stabilizer.*, gyro.*, acc.*,
%   indi.a_res_*, indi.e_r_*, indi.tau_*, indi.alp_*, rnn.pred_*, rnn.clamped, ctrltarget.*,
%   motor.m*, motor.m*_rpm, rpm.m*, rpmf.m*, pm.vbat, locSrv.*, diagnostic groups, usd.runTag).

Two of these deserve comment. The measured residual is recorded \emph{in every controller mode},
including those that make no use of it. Training data is collected under the uncompensated baseline.
The network prediction is recorded \emph{whether it is enabled or not}. Comparing prediction to
measurement requires both on the same flight without the prediction driving the loop.

The logging configuration is fixed before data collection and shared by the collection and
evaluation flights. A configuration change between the two would silently make the two halves of
the dataset incomparable.

\subsection{Time alignment across vehicles}

Each vehicle timestamps its samples against its own boot, so the logs share no clock. Recording is
started by a single broadcast, which aligns the vehicles to within the broadcast jitter plus one
logging period. Residual error comes from propagation, task wake phase and crystal drift between
vehicles; the last is a rate error and accumulates over a flight.

% CHANGED (V6): the alignment method as implemented in flying_drone_stack/tools/merge_usd_logs.py
%   (module docstring L1-32; --meta/--roles alignment against each vehicle's own commanded trajectory,
%   L339-342, L393-445; radio cross-correlation with offset and drift, L80-170, as a cross-check).
% DECISION D-05 (H9): the last sentence of the next paragraph ("to be confirmed against real logs on
%   the first two-vehicle flight") is project status. Recommended: "The measured offset and drift are
%   reported for every two-vehicle session." Text left unchanged.
The merge tool therefore \emph{measures} the alignment rather than assuming it. Each vehicle's
onboard position trace is aligned against that vehicle's own commanded trajectory, which is known
from the scenario definition and differs between vehicles, and the quality of the fit is reported.
Where radio logs of the same flight exist, which share a common clock, cross-correlating each
onboard trace against its radio trace additionally gives the offset and the drift as a
cross-check. The predicted alignment is a few milliseconds; that figure is to be confirmed
against real logs on the first two-vehicle flight.

% -------------------------------------------------------------------------------------
\section{Protocol}
\label{sec:setup-protocol}

\subsection{Hardware gate}
\label{sec:setup-gate}

No data collected before the following sequence passes is used. The gate is not an experiment; it
is the check that nothing is broken, so that what is measured afterwards means something.

\begin{enumerate}
  \item \textbf{Validate.} Single vehicle, geometric first and then INDI, on a sequence of trajectories from hover to the
        formation trajectories.
        Confirm the gains read back as configured.
  \item \textbf{Tune, then freeze.} Any remaining attitude tuning happens here and only here. The
        gains are then frozen for the remainder of the campaign.
  \item \textbf{Two vehicles at large separation.} Confirm $a_{\mathrm{res}}$ is non-zero and
        correctly signed; confirm the measured time alignment.
\end{enumerate}

The order carries the reasoning. Tuning before validating means tuning against an unknown fault.
Flying two vehicles before the freeze means the gains move underneath the dataset.

Stage~1 cannot confirm that the compensation terms are \emph{correct}, only that they are not
catastrophic: with no neighbour present $a_{\mathrm{res}} \approx 0$, so a compensation term is
dormant and a clean single-vehicle flight demonstrates nothing about it. Stage~3 is the first point
at which the sign of the compensation is observable at all, which is why it is part of the gate
rather than the first experiment.
% DECISION D-01: Stage 3 "confirm a_res is ... correctly signed" presupposes the sign decision
%   (Strategy 1's configured res_sign = +1 vs the subtraction of Ch3 eq. model-sign).
%   CLOSING PASS (D-11): the values above for Strategy 1 are those of the 'ours' block (controller_step INDI,
%   ctrl_mode 3). The frozen S1 block of crazyswarm2/crazyflie/config/study_strategies.json selects a
%   different INDI controller (stabilizer.controller 9, controller_omar_indi.c), which subtracts the
%   residual (L260, L268) with matched filters and runs a vertical integral (Kpos_Iz 1.5). Which one
%   Strategy 1 is, is decision D-11.

\subsection{Data collection}
\label{sec:setup-collect}

Collection flights are flown under the \emph{uncompensated} baseline of
Section~\ref{sec:ctrl-s0}. This is not a convenience: a compensating controller partly cancels the
disturbance, so the residual it records is the one that survived compensation, which is not the
quantity a predictive model should be fitted to.
% DECISION D-04: the reason given contradicts Ch3 Sec. 3.3.1 (fourth paragraph) and Shi et al.
%   2020 Sec. III-C. Reasons that hold are listed in author_decisions.md. Text left unchanged.

Separations are approached from large to small, and a tighter configuration is attempted only after
the looser ones are clean.

\textbf{The set must span both vertical and lateral excitation} --- A1/A2/A3 \emph{and} A4/A5/A7
--- for the reason given in Section~\ref{sec:setup-scenarios}. A collection campaign that omits the
laterally excited scenarios cannot support the predictive strategies, and the omission is not
recoverable without re-flying.

\subsection{Comparison campaign}
\label{sec:setup-campaign}

% CHANGED (P5-23): the number of repetitions is an open protocol parameter.
Each strategy flies the frozen library. Because the control law is selected at runtime
(Section~\ref{sec:ctrl-runtime}), strategies are interleaved within a session rather than flown on
separate days, which pairs them against the same battery state, temperature and calibration.
Each condition is repeated to allow a distribution rather than a single number to be reported.
\todo{Fix the number of repetitions per (scenario, strategy) cell before the campaign.}

Each strategy is compared under identical trajectories, formations, hardware and separations, with
its own gains frozen after the hardware gate, not identical control structure — this holds
uniformly across all of them, not just a subset.

% NEW (from Ch4 Sec. 4.4): the two kinds of flight Strategy 2 needs.
For Strategy~2 two kinds of flight are distinguished. On \emph{prediction-off} flights, flown
under Strategy~0 with the trained weights loaded, the prediction is recorded but not used, so
predicted and measured residuals can be compared without the prediction having moved the vehicle.
On \emph{prediction-on} flights the prediction drives the controller, and these flights are the
ones compared with the other strategies.

\subsection{Metrics}
\label{sec:setup-metrics}

Three axes, corresponding to the research questions.

% NEW (P5-11): definitions, so that every metric is reproducible from the recorded variables.
\paragraph{Notation and windows.} For vehicle $i$, $p_{\mathrm{cmd},i}(t)$ is the commanded position
recorded onboard and $p_i(t)$ the measured position, and $e_{p,i}(t) = p_{\mathrm{cmd},i}(t) - p_i(t)$ is its
position error, with the sign of $e_p$ in Chapter~\ref{ch:control}; $e_1, e_2, e_3$ are the unit
vectors of the world frame.
% CHANGED (closing pass, diff audit): e_i renamed e_{p,i} (it collided with the unit vectors e_l)
%   and its sign aligned with e_p = p_des - p of Chapter 4.
Each scenario defines an analysis window $W$, the set of samples in which its commanded geometry is
in its nominal configuration, excluding take-off, landing and the settling intervals of the
scenario definition. For the swap scenarios (A8, B3) an event window
$W_c = [t_c - \Delta, t_c + \Delta]$ is centred on each crossing time $t_c$, defined as the instant
of minimum \emph{commanded} horizontal distance between the two vehicles concerned; because it is
taken from the commanded trajectory, it does not depend on how a controller responded.
\todo{Confirm before the campaign: the source of $p_i$ (onboard estimate or motion-capture pose,
both recorded), the window boundaries per scenario, and the half-width $\Delta$ of the event window.}
For a quantity $q$ sampled at instants $t_k \in W$, $\operatorname{RMS}_W(q) =
\big(\tfrac{1}{\lvert W\rvert}\sum_{k} \lVert q(t_k)\rVert^2\big)^{1/2}$.

\paragraph{Tracking error.} Position RMSE, total and per axis, and the maximum deviation from the
commanded trajectory: $\operatorname{RMS}_W(e_{p,i})$, $\operatorname{RMS}_W(e_{p,i}^\top e_\ell)$ for
$\ell = 1,2,3$, and $\max_{t\in W}\lVert e_{p,i}(t)\rVert$, with the vertical maximum
$\max_{t\in W} \lvert e_{p,i}(t)^\top e_3\rvert$ reported separately because downwash acts mainly
vertically. Alongside these, the attitude tracking error, reported as the RMSE and peak
of the geometric attitude error $e_R$. A compensator that improves position by working the attitude
loop harder has not improved tracking so much as moved the cost, and only the attitude figure makes
that visible.
% CHANGED (P5-04): e_R is in the logging configuration.
The logging configuration of Section~\ref{sec:setup-recorded} records $e_R$ and its norm under every
controller mode, so $\operatorname{RMS}_W(e_R)$ and $\max_{t\in W}\lVert e_R(t)\rVert$ are computed
directly.
% CHANGED: code, usd_study_config.txt (indi.e_r_x/y/z, indi.e_r_norm); lib.rs:controller_step
%   (written before the geometric/INDI branch, L1972-1976).
% REMOVED (P5-04): "The logging configuration described in Section~\ref{sec:setup-recorded} records
%   attitude but not $e_R$ itself; that signal will be added before the comparison campaign."

\paragraph{Residual rejection.} The error in the \emph{separation between vehicles}, measured
against the commanded separation. This is a separation error, not the residual itself; the residual
that remains is reported next to it. Downwash acts as a bias, the lower vehicle sags, widening the
vertical gap, so the mean error is reported as the primary quantity and the spread as secondary. A
compensator that works moves the mean back toward the commanded value. Reporting only the spread
would hide precisely the effect being studied.
% CHANGED (closing pass, examiner): "closing the gap" -> "widening the vertical gap": the downwash pushes
%   the LOWER vehicle down, so the vertical separation grows (consistent with e_sep below, which is then
%   positive).
% NEW (P5-11): the separation error, defined.
For an upper vehicle $u$ and a lower vehicle $b$ in a stack,
$e_{\mathrm{sep}}(t) = (p_u(t) - p_b(t))^\top e_3 - (p_{\mathrm{cmd},u}(t) - p_{\mathrm{cmd},b}(t))^\top e_3$;
the reported quantities are its mean and standard deviation over $W$ (or $W_c$).

Alongside it, the magnitude of the measured residual that remains after compensation, and for every
strategy carrying a learned model, the prediction against the measurement.
% NEW (P5-11): definitions for the remaining residual, the prediction error and the transfer penalty.
The remaining residual is $\operatorname{RMS}_W(a_{\mathrm{res}})$, reported with its vertical
component. The prediction error of a model is evaluated on prediction-off flights as
$\epsilon_{\mathrm{pred}} = \operatorname{RMS}_W(\hat a_{\mathrm{res},z} - a_{\mathrm{res},z})$ together
with its mean, which separates a bias from scatter. The transfer penalty that RQ3 asks for is the
difference between $\epsilon_{\mathrm{pred}}$ on the three-vehicle scenarios and on the two-vehicle
scenarios, for the same weights, absolute and relative to the two-vehicle value.
\todo{Fix which two-vehicle scenarios are paired with each three-vehicle scenario for the transfer
penalty (for example B1 with A2 at matching height differences).}

\paragraph{Robustness and cost.} Control effort; behaviour as separation tightens; abort and
failure rate; onboard computation per control step; sensing prerequisites; and the quantity of
training data required to reach a given prediction accuracy.
% NEW (P5-11): an effort definition from recorded variables.
Control effort is reported as $\operatorname{RMS}_W$ of the commanded torque and of the change of
the motor commands between samples, and as the fraction of samples in which any motor command is
at its limit.
\todo{Confirm the effort measures before the campaign.}
% CHANGED (closing pass, diff audit): the effort definition is a proposal, so it is put under a todo.
% DECISION D-09: whether the training-data quantity is measured by a data-size sweep (as RQ4 implies)
%   and which hypotheses are stated with their falsifiers; see author_decisions.md.

% NEW (P5-12): how repetitions are aggregated.
\subsection{Statistical treatment}
\label{sec:setup-stats}

Each metric is computed per flight, then summarised per (scenario, strategy) cell over the
repetitions. Because strategies are interleaved within sessions, the comparison of a strategy with
Strategy~0 can be made on pairs of flights from the same session. The null scenarios bound the
measurement floor (Section~\ref{sec:setup-null}) against which a difference between strategies is
judged.
\todo{Fix before the campaign: the summary statistic per cell, the paired test or interval method,
the effect-size measure, and the rule by which a difference is compared with the measurement floor
of the null scenarios.}
% CHANGED (closing pass, diff audit): the earlier sentence stated a decision rule ("smaller than the
%   spread on C1-C3 is indistinguishable") that no source or protocol document fixes; it is now part
%   of the todo, and pairing is "can be made", not a rule.

% -------------------------------------------------------------------------------------
\section{The role of simulation}
\label{sec:setup-sim}

% CHANGED (P5-16): the interaction model is named; tense and findings per H1.
A software-in-the-loop simulator runs the same controller source that flies, with the published
Neural-Swarm2 interaction model~\cite{shi2022neuralswarm2} providing the interaction force. Its role
in this thesis is deliberately limited.
% VERIFY (V7, partly closed): architecture identical to the flight network (neuralswarm.py L24-56:
%   phi 6->25->40->40->20, rho 20->40->40->40->1, H = 20; gate L78 as residual_nn.rs). The weight
%   files (backend/data/neuralswarm2/*.pth) were added to upstream Crazyswarm2 by W. Hoenig, a
%   Neural-Swarm2 author (crazyswarm2 commit 51a12dd, 2024-01-10). Open: whether they are the networks
%   of the paper -- ask the supervisor.

\textbf{Simulation is used to disprove implementations, never to validate a method.} It is used to
check that every scenario in the library executes and produces the commanded geometry under each
implemented controller, and to find defects in flight code that single-vehicle testing cannot
reveal, before the vehicles fly. It cannot support a claim about how well a method works, for a
reason specific to this subject: the interaction force in simulation is produced by a learned model
of the same family as the one being fitted --- in fact of the same architecture --- and learning a
network's output with another network is a substantially easier problem than learning real air.
% REMOVED (P5-06, H1): "It was used throughout development, and its role in this thesis is
%   deliberately limited." and "It established that every scenario in the library executes and
%   produces the commanded geometry under both controllers, and it exposed two defects in flight
%   code that single-vehicle testing could not have found." (project simulation findings)

No result in the two-robot results chapter or~the multi-robot results chapter rests on simulation.

% -------------------------------------------------------------------------------------
\section{Threats to validity}
\label{sec:setup-validity}

% CHANGED (P5-17): internal vocabulary in the "radio vs onboard clock" row.
\begin{table}[t]
  \centering
  \caption{Threats to validity and the design decisions taken against them.}
  \label{tab:setup-threats}
  \small
  \begin{tabular}{@{}p{0.28\linewidth}p{0.34\linewidth}p{0.30\linewidth}@{}}
    \toprule
    Threat & Mitigation & Residual risk \\
    \midrule
    Tuning confound & Each strategy tuned on the hardware gate; those gains frozen before comparison flights; runtime selection allows interleaving & Per-strategy tuning effort is in scope for RQ4 (§\ref{sec:ctrl-fairness}) \\
    Measurement floor & Null scenarios C1--C3 bound it under the same pipeline & --- \\
    Model error read as interaction & Null scenarios detect it; airframe constants identified once and frozen & Battery discharge within a flight is not compensated \\
    Training/test leakage & Held-out scenarios; split by contiguous flight segment, never by sample & --- \\
    Session effects & Strategies interleaved within a session; repeated conditions & --- \\
    Silent loss of the residual & $a_{\mathrm{res}} \neq 0$ verified at the start of every session & --- \\
    Radio vs onboard clock & Dataset recorded onboard at \SI{500}{\hertz}; analysis uses the merged onboard logs and the commanded setpoint recorded onboard & Residual-time alignment across vehicles depends on merge quality; the merge refuses an alignment whose position fit to the vehicle's own commanded trajectory exceeds \SI{15}{\centi\metre} RMS, unless overridden explicitly \\
    Mocap identity / marker swap & Pre-takeoff check of the onboard estimate against motion capture; role verification from merge & Residual risk if a session bypasses the gate \\
    Simulation-to-hardware gap & No claim rests on simulation & --- \\
    Single airframe & --- & \textbf{Not mitigated.} Generalisation beyond this platform is discussed, not demonstrated \\
    \bottomrule
  \end{tabular}
\end{table}
% CHANGED (V8): merge_usd_logs.py L432-438 (refuses at RMS >= 15 cm; override flag
%   --tag-quality-only, L350-353). A fit to the commanded trajectory also contains the tracking
%   error itself, which is why an override exists for close stacks.

% CHANGED: precedent for the split, from a held source (R02 Sec. VII-C).
Splitting the training and test data by contiguous flight segment rather than by sample deserves a
note, because the obvious alternative is badly wrong here. At \SI{500}{\hertz} consecutive samples
are nearly identical, so a random per-sample split places a near-duplicate of almost every test
point in the training set and reports a validation error that means nothing. Neural-Swarm2 splits
its data the same way, assigning pieces of about \SI{20}{\second} to training or validation rather
than individual samples~\cite{shi2022neuralswarm2}.

Two threats are not mitigated and are stated as limitations rather than managed. The results
concern a single airframe; quantities are normalised where a normalisation is defensible, but
generalisation to other platforms is argued rather than shown. And battery discharge changes the
thrust constant slightly within a flight, which appears as a slow drift in the measured residual;
the effect is bounded by the null scenarios but not removed.

% -------------------------------------------------------------------------------------
\section{Summary}
\label{sec:setup-summary}

% CHANGED (P5-21): numbers from the code.
Two or three identical brushless quadrotors, each carrying a rotor-speed source and an SD card,
flying sixteen frozen scenarios inside a tracked volume of roughly \SI{2}{\metre} $\times$
\SI{4}{\metre} $\times$ \SI{1.7}{\metre}, with a control loop at \SI{1000}{\hertz} and \num{96}
variables recorded onboard at \SI{500}{\hertz} independently of team size. Each strategy is tuned
on the hardware gate. Those
gains are frozen before the comparison and training-collection flights; strategies are selected at
runtime and interleaved within sessions; and three null scenarios
establish what a measurement of zero looks like under the same pipeline.

% CHANGED (P5-22): internal workflow label removed from the todo.
\todo{Continue revising to past tense for each scenario and controller row as the comparison
campaign's repetitions are completed. Record any deviation from a per-strategy gain block, frozen
after the hardware gate, or from scenario parameters in this chapter.}
